\chapter{Why an artefact substitutes for a title}
\label{ch:artefact}

\section*{What this chapter is for}

Your curriculum vitae does not say \emph{AI Engineer}. Somebody has to decide,
from the outside, whether you can do the work. This chapter is about what
actually closes that gap, and \dash{} more importantly \dash{} about what it
does not close, because mistaking the second for the first is how a month gets
spent in the wrong place.

\section{The mechanism}

A hiring decision at this level is a decision under uncertainty about future
behaviour. The evidence available is: what you have done before, what you say
about it, and what somebody can see for themselves.

A job title is a compressed, second-hand summary of the first. It is
convenient, it is what filters are built on, and it is weak evidence: it tells a
reader that an organisation once applied a label, not what you did under it.

An artefact is the third kind. It is the only evidence in the set that the
reader can inspect without trusting you. \reported{Deployed, clickable systems
are reported to beat notebooks, and some companies ask for a portfolio up front
specifically to separate people who have built something from people who have
followed a tutorial.}{field-guide-get-hired}

\begin{kitnote}
This is why the substitution works, and it is worth being precise about it:
an artefact does not \emph{prove} more than a title. It requires less trust. A
reader deciding under uncertainty prefers the evidence they can check to the
evidence they must accept, and that preference is the whole of the mechanism.
\end{kitnote}

\section{What it substitutes for, exactly}

\begin{kittable}{@{}lXl@{}}
\textbf{The doubt} & \textbf{Stated as a question} & \textbf{Artefact?} \\
\midrule
Domain & Has this person built one of these? & yes \\
Currency & Have they built one recently? & yes \\
Depth & Do they understand it or did they assemble it? & partly \\
Judgement & Do they know what not to build? & partly \\
Scope & Have they owned something large and ambiguous? & no \\
Tenure & Will they still be here in two years? & no \\
Collaboration & What are they like when someone disagrees? & no \\
\end{kittable}

The first two rows are where artefacts are decisive, and they are exactly the
two rows a missing title creates. That is the good news and it is genuinely
good: the specific doubt your background raises is the specific doubt an
artefact removes.

The last three rows are the ones to be honest about. \judgement{An artefact is
evidence about a thing you built. It is not evidence about how you behave over
two years, how you behave when somebody senior disagrees with you, or whether
you can hold an unscoped problem for six months. Those are assessed elsewhere
in the loop, from a different body of evidence, and no amount of building moves
them.} Chapter~\ref{ch:classes} is about that division and about checking the
immovable ones first.

The two \emph{partly} rows are the interesting ones, and they are why
Chapter~\ref{ch:legible} exists. Whether an artefact demonstrates depth and
judgement, rather than assembly, is a property of how it is built and written
up, not of what it does.

\section{Why this is stronger at Staff than at Senior}

It would be reasonable to expect the opposite \dash{} that a more senior role
cares less about a side project.

\judgement{The reverse holds, for a reason that comes straight from
Chapter~\ref{ch:staff}.} A Senior role is hiring for execution, and execution
is what a job history evidences well. A Staff role is hiring for judgement
about unscoped problems, and a job history evidences judgement poorly: it tells
the reader what got built, not which of the alternatives were rejected and why.

An artefact built on your own time is an unscoped problem you scoped yourself.
Every decision in it was yours, including the decision about what not to do.
That is the exact evidence the level is short of, and it is why a small,
well-argued artefact can outweigh a large vague one.

\section{The honest limit, stated once}

An artefact substitutes for a \emph{title}. It does not substitute for
\emph{tenure}, and it does not substitute for demonstrated scope in a real
organisation.

\begin{kittrap}
It is tempting to believe that a sufficiently good portfolio answers every
doubt, because the portfolio is the part you control and the part that responds
to effort. It answers two doubts extremely well and several others not at all.

The failure mode this produces is specific: a candidate who keeps building,
because building is legible and feels like progress, while the thing that is
actually blocking them sits in a class no amount of building reaches. Working
out which class you are in is Chapter~\ref{ch:classes}, and it is worth an
afternoon before it is worth another month of evenings.
\end{kittrap}

\begin{drill}
Write the two sentences an interviewer would use to argue against hiring you,
in their words, not yours. Then, for each, say whether an artefact could answer
it. If both could be answered by building something, Part~\ref{part:evidence}
is your next month. If neither could, it is not \dash{} and you have saved the
month.
\end{drill}

\section*{What this chapter does not claim}

It does not claim artefacts are necessary. People move into this work without
them, through internal transfers and through employers who hire on trajectory.
The claim is narrower: where a title is missing and the reader has no other
basis for trust, a thing they can inspect is the cheapest way to supply one.

It does not claim any particular artefact will have this effect. What is
reported is that portfolio evidence is used, in some loops, as a filter with
that intent. How often, and with what weight, is not something this book knows.
